\chapter{Complexity — Circuit Sat Reduction Machine}
%%% generated by the blueprint pipeline from TCSlib.Complexity.CircuitComplexity.CircuitSatReductionMachine
\section{Overview}

Declarations of the Lean module \texttt{TCSlib.Complexity.CircuitComplexity.CircuitSatReductionMachine}.

\section{Declarations}

\begin{definition}[Unary tapes from three lengths]
\lean{BoolCircuit.CktSatReduction.Emitter.tapesOf}
\label{BoolCircuit.CktSatReduction.Emitter.tapesOf}
\leanok
\uses{Turing.UnaryTape.ones}
Given three natural numbers $v_0,v_1,v_2$, form a family of three tapes by assigning to
tape $i$ the unary tape of length $v_i$.
\end{definition}

\begin{definition}[Head positions from three lengths]
\lean{BoolCircuit.CktSatReduction.Emitter.headsOf}
\label{BoolCircuit.CktSatReduction.Emitter.headsOf}
\leanok
Given three natural numbers $v_0,v_1,v_2$, assign to each tape $i$ the integer head
position equal to $v_i$, viewed as an integer.
\end{definition}

\begin{theorem}[Updating one unary tape]
\lean{BoolCircuit.CktSatReduction.Emitter.tapesOf_update}
\label{BoolCircuit.CktSatReduction.Emitter.tapesOf_update}
\leanok
\uses{BoolCircuit.CktSatReduction.Emitter.tapesOf, Turing.UnaryTape.ones}
\difficulty{4}
Let $v$ assign a natural number to each of three tape indices, and let $t$ be one of
those indices and $m$ a natural number. Replacing tape $t$ by the unary tape of length
$m$ in the family of tapes determined by $v$ gives exactly the family determined by
replacing $v(t)$ with $m$.
\end{theorem}

\begin{theorem}[Updating work-tape head positions]
\lean{BoolCircuit.CktSatReduction.Emitter.headsOf_update}
\label{BoolCircuit.CktSatReduction.Emitter.headsOf_update}
\leanok
\uses{BoolCircuit.CktSatReduction.Emitter.headsOf}
\difficulty{4}
For any function $v$ assigning a natural number to each of three tape indices, any index
$t\in\{0,1,2\}$, and any natural number $m$, converting the values of $v$ to integers
and then replacing the value at $t$ by $m$ gives the same function as first replacing
$v(t)$ by $m$ and then converting all values to integers.
\end{theorem}

\begin{theorem}[Zero update of head positions]
\lean{BoolCircuit.CktSatReduction.Emitter.headsOf_update_zero}
\label{BoolCircuit.CktSatReduction.Emitter.headsOf_update_zero}
\leanok
\uses{BoolCircuit.CktSatReduction.Emitter.headsOf, BoolCircuit.CktSatReduction.Emitter.headsOf_update}
\difficulty{1}
For any function $v$ from the three tape indices to the natural numbers and any index
$t\in\{0,1,2\}$, replacing the integer-valued version of $v$ at $t$ by zero gives the
same function as first replacing $v(t)$ by zero and then converting all values to
integers.
\end{theorem}

\begin{definition}[Shifted input position]
\lean{BoolCircuit.CktSatReduction.Emitter.ip}
\label{BoolCircuit.CktSatReduction.Emitter.ip}
\leanok
For a Boolean list $w$ and a natural number $i$ with $i\\leq |w|$, let the corresponding
input-head position be $i+1$ in the range of $|w|+2$ positions.
\end{definition}

\begin{theorem}[Rightward input-head movement]
\lean{BoolCircuit.CktSatReduction.Emitter.move_ip}
\label{BoolCircuit.CktSatReduction.Emitter.move_ip}
\leanok
\uses{BoolCircuit.CktSatReduction.Emitter.ip, Turing.moveInputPos, Turing.moveInputPos_pos_of_ne_right}
\difficulty{3}
Let $w$ be a Boolean input and let $i<|w|$. If the input head is at position $i+1$ in
the indexing that places the two boundary cells at positions $0$ and $|w|+1$, then a
rightward move places it at position $i+2$.
\end{theorem}

\begin{theorem}[Machine execution for an emitter instruction suffix]
\lean{BoolCircuit.CktSatReduction.Emitter.em_loop}
\label{BoolCircuit.CktSatReduction.Emitter.em_loop}
\leanok
\uses{BoolCircuit.CktSatReduction.Emitter.Instr, BoolCircuit.CktSatReduction.Emitter.St, BoolCircuit.CktSatReduction.Emitter.apply_act, BoolCircuit.CktSatReduction.Emitter.apply_ctl, BoolCircuit.CktSatReduction.Emitter.back_walk, BoolCircuit.CktSatReduction.Emitter.emTM, BoolCircuit.CktSatReduction.Emitter.exec, BoolCircuit.CktSatReduction.Emitter.headsOf, BoolCircuit.CktSatReduction.Emitter.length_prog_le, BoolCircuit.CktSatReduction.Emitter.print_walk, BoolCircuit.CktSatReduction.Emitter.prog, BoolCircuit.CktSatReduction.Emitter.step_some, BoolCircuit.CktSatReduction.Emitter.tapesOf, BoolCircuit.CktSatReduction.Emitter.tr, BoolCircuit.GateKind, Turing.Cfg, Turing.MultiTapeTM.runFrom, Turing.MultiTapeTM.runFrom_add, Turing.MultiTapeTM.runFrom_succ_eq_step, Turing.MultiTapeTM.runFrom_zero, Turing.UnaryTape.wrT, Turing.moveInputPos_zero}
\difficulty{7}
Let $w$ be a Boolean list, let $k$ be a gate kind, let $c$ be one of $0,1,2$, and let
$p$ be an input-head position among the $|w|+2$ positions. Let $v$ assign a natural
number to each of three work tapes, and suppose every value of $v$ is at most $M$. For
any instruction list $r$ and program counter $pc<64$ such that the program for $(k,c)$
starting at $pc$ is exactly $r$, and any output list $o$, there are a time $T\le
|r|(2M+3)$ and a program counter $p_e<64$ at which the program has ended such that,
after $T$ steps from the configuration in emission state $(k,c,pc)$ with input-head
position $p$, unary tapes and heads specified by $v$, and output $o$, the machine is in
emission state $(k,c,p_e)$ with the same input position, tapes, and heads, and output
$o$ followed by the execution of $r$ under $v$.
\end{theorem}

\begin{definition}[Emitter machine configuration]
\lean{BoolCircuit.CktSatReduction.Emitter.cfgS}
\label{BoolCircuit.CktSatReduction.Emitter.cfgS}
\leanok
\uses{BoolCircuit.CktSatReduction.Emitter.ES, BoolCircuit.CktSatReduction.Emitter.ES.vec, BoolCircuit.CktSatReduction.Emitter.St, BoolCircuit.CktSatReduction.Emitter.headsOf, BoolCircuit.CktSatReduction.Emitter.ip, BoolCircuit.CktSatReduction.Emitter.tapesOf, Turing.Cfg}
For a Boolean input list $w$, an emitter state $s$, an index $i$ at most the length of
$w$, and an output list, form a configuration whose machine state is the read phase
specified by $s$, whose input-head position is the shifted position $i+1$, and whose
three work tapes and their head positions are given by the components of $s$. Its output
tape contains the supplied output list.
\end{definition}

\begin{theorem}[One read-state transition]
\lean{BoolCircuit.CktSatReduction.Emitter.rd_step}
\label{BoolCircuit.CktSatReduction.Emitter.rd_step}
\leanok
\uses{BoolCircuit.CktSatReduction.Emitter.ES, BoolCircuit.CktSatReduction.Emitter.ES.vec, BoolCircuit.CktSatReduction.Emitter.cfgS, BoolCircuit.CktSatReduction.Emitter.emTM, BoolCircuit.CktSatReduction.Emitter.headsOf, BoolCircuit.CktSatReduction.Emitter.step_some, BoolCircuit.CktSatReduction.Emitter.tapesOf, BoolCircuit.CktSatReduction.Emitter.tr, Turing.Action.apply, Turing.FinTM.inputSymbol_at, Turing.MultiTapeTM.runFrom, Turing.MultiTapeTM.runFrom_succ_eq_step, Turing.MultiTapeTM.runFrom_zero}
\difficulty{3}
Let $w$ be a Boolean input, $s$ an emitter state, $i\le |w|$, and $out$ an output list.
Starting the emitter machine in the read state with phase $s.\varphi$, input-head
position corresponding to $i$, unary work tapes encoding the coordinates of $s$, and
output $out$, its configuration after one step is exactly the result of applying the
action prescribed by the transition rule for that read state, the input symbol at
position $i$ (if present), and the symbols under those work-tape heads.
\end{theorem}

\begin{theorem}[Gate emission runtime]
\lean{BoolCircuit.CktSatReduction.Emitter.gate_run}
\label{BoolCircuit.CktSatReduction.Emitter.gate_run}
\leanok
\uses{BoolCircuit.CktSatReduction.Emitter.ES, BoolCircuit.CktSatReduction.Emitter.ES.vec, BoolCircuit.CktSatReduction.Emitter.St, BoolCircuit.CktSatReduction.Emitter.apply_act, BoolCircuit.CktSatReduction.Emitter.cfgS, BoolCircuit.CktSatReduction.Emitter.clear_walk, BoolCircuit.CktSatReduction.Emitter.emTM, BoolCircuit.CktSatReduction.Emitter.em_loop, BoolCircuit.CktSatReduction.Emitter.emitBits, BoolCircuit.CktSatReduction.Emitter.exec, BoolCircuit.CktSatReduction.Emitter.headsOf, BoolCircuit.CktSatReduction.Emitter.headsOf_update_zero, BoolCircuit.CktSatReduction.Emitter.ip, BoolCircuit.CktSatReduction.Emitter.length_prog_le, BoolCircuit.CktSatReduction.Emitter.prog, BoolCircuit.CktSatReduction.Emitter.step_some, BoolCircuit.CktSatReduction.Emitter.tapesOf, BoolCircuit.CktSatReduction.Emitter.tapesOf_update, BoolCircuit.CktSatReduction.Emitter.tr, BoolCircuit.GateKind, Turing.Cfg, Turing.MultiTapeTM.runFrom, Turing.MultiTapeTM.runFrom_add, Turing.MultiTapeTM.runFrom_succ_eq_step, Turing.MultiTapeTM.runFrom_zero, Turing.UnaryTape.ones, Turing.UnaryTape.ones_zero, Turing.UnaryTape.update_ones_succ, Turing.UnaryTape.wrT, Turing.moveInputPos_zero}
\difficulty{6}
Let $w$ be a finite Boolean input, let $s=(\phi,z,a,b)$ be an emitter state, let
$\kappa$ be a gate kind, let $c\in\{0,1,2\}$ be a gate code, and let $i\le |w|$. For any
natural number $M$ such that $z,a,b\le M$, and any output string $out$, there is a
number of steps $T\le 200(M+1)$ after which the emitter machine, started in its
gate-emission state with program counter zero, input-head position corresponding to $i$,
unary work tapes and work-head positions encoding $(z,a,b)$, and output $out$, is in its
read state with phase $g$, vertex counter $z+1$, both argument counters zero, the same
input-head position, unary work tapes and work-head positions encoding $(z+1,0,0)$, and
output equal to $out$ followed by the bits emitted for gate kind $\kappa$, code $c$, and
counters $(z,a,b)$.
\end{theorem}

\begin{theorem}[Machine execution for an emitter run]
\lean{BoolCircuit.CktSatReduction.Emitter.em_main_run}
\label{BoolCircuit.CktSatReduction.Emitter.em_main_run}
\leanok
\uses{BoolCircuit.CktSatReduction.Emitter.ES, BoolCircuit.CktSatReduction.Emitter.ES.incTo, BoolCircuit.CktSatReduction.Emitter.ES.vec, BoolCircuit.CktSatReduction.Emitter.ES.vec_incTo, BoolCircuit.CktSatReduction.Emitter.apply_act, BoolCircuit.CktSatReduction.Emitter.apply_ctl, BoolCircuit.CktSatReduction.Emitter.cfgS, BoolCircuit.CktSatReduction.Emitter.ctl, BoolCircuit.CktSatReduction.Emitter.emRun, BoolCircuit.CktSatReduction.Emitter.emTM, BoolCircuit.CktSatReduction.Emitter.emitBits, BoolCircuit.CktSatReduction.Emitter.gate_run, BoolCircuit.CktSatReduction.Emitter.headsOf_update, BoolCircuit.CktSatReduction.Emitter.lstep, BoolCircuit.CktSatReduction.Emitter.move_ip, BoolCircuit.CktSatReduction.Emitter.rd_step, BoolCircuit.CktSatReduction.Emitter.tapesOf_update, BoolCircuit.CktSatReduction.Emitter.tr, Turing.MultiTapeTM.runFrom, Turing.MultiTapeTM.runFrom_add, Turing.UnaryTape.ones, Turing.UnaryTape.update_ones_succ, Turing.UnaryTape.wrT}
\difficulty{6}
For every Boolean input list $w$, every Boolean list $l$, emitter state $s$, index $i\le
|w|$, and output list $o$, suppose dropping the first $i$ bits of $w$ gives $l$ and each
of the three counters in $s$ is at most $i$. Then there is a time $T\le
(|l|+1)\,200(|w|+1)$ such that the machine started in the configuration determined by
$w,s,i,o$ has halted after $T$ steps and its output is $o$ followed by the result of
running the emitter from $s$ on $l$.
\end{theorem}

\begin{theorem}[Running time of the gate clause machine]
\lean{BoolCircuit.CktSatReduction.Emitter.emTM_computes}
\label{BoolCircuit.CktSatReduction.Emitter.emTM_computes}
\leanok
\uses{BoolCircuit.CktSatReduction.Emitter.ES.vec, BoolCircuit.CktSatReduction.Emitter.cfgS, BoolCircuit.CktSatReduction.Emitter.emTM, BoolCircuit.CktSatReduction.Emitter.em_main_run, BoolCircuit.CktSatReduction.Emitter.gateClauses, BoolCircuit.CktSatReduction.Emitter.headsOf, BoolCircuit.CktSatReduction.Emitter.tapesOf, Turing.Cfg.init, Turing.FinTM.ComputesFunInTime, Turing.FinTM.computesInTime_iff, Turing.MultiTapeTM.initCfg, Turing.MultiTapeTM.runFrom_add, Turing.MultiTapeTM.runFrom_of_halt, Turing.UnaryTape.ones_zero}
\difficulty{4}
There is a finite-state multi-tape Turing machine that, on every Boolean input of length
$n$, halts with output equal to the gate-clause encoding of that input and uses exactly
$200(n+1)^2$ units of time.
\end{theorem}
